\documentclass{article}

\usepackage{fullpage}
\usepackage{url}
\usepackage{graphicx}
\usepackage{amsmath}
\usepackage{physics}
\usepackage{mathtools}
\usepackage{bbold}

\usepackage{lmodern}
\usepackage[T1]{fontenc}
\usepackage[fontsize=13.5pt]{fontsize}

\DeclareMathOperator{\erfi}{erfi}
\newcommand{\R}{\ensuremath{\mathbb{R}}}
\newcommand{\C}{\ensuremath{\mathbb{C}}}

\title{Complex Systems Excercise 4}
\author{Idan Stark}
\date{\today}

\begin{document}

\maketitle

\section{Fokker - Planck Equation for an Ornstein - Uhlenbeck Process}
In this section we start from the Langevin equation for the Ornstein - Uhlenbeck Process:
\begin{equation}
    \frac{dv}{dt} = -\frac{v}{\tau} + \xi
\end{equation}
With
\begin{equation}
    \langle \xi(t) \xi(t') \rangle = A \delta(t - t')
\end{equation}
\subsection{The Fokker - Planck}
Using the general form of a Langevin equation:
\begin{equation}
    \frac{dv}{dt} = -\mu F(v) + \xi
\end{equation}
We can write:
\begin{equation}
    F = v \qquad \mu = \tau^{-1}
\end{equation}
And so the Fokker-Planck equation comes from the continuity equation, where the probability density current is a result of the force and of diffusion:
\begin{equation}
    \frac{dp}{dt} = \mu \partial_v \left[F(v) p\right] + \partial_v \left[D \partial_v p\right]
\end{equation}
We can set the value for $D$ based on the $\mu = 0$ limit. At that limit, for the delta noise:
\begin{equation}
    D = A/2
\end{equation}
Plugging this, and the values for $\mu$ and $F(v)$:
\begin{equation}
    \frac{dp}{dt} = \frac{1}{\tau} \partial_v \left[pv\right] + \partial_v \left[\frac{A}{2} \partial_v p\right]
\end{equation}
We get the Fokker - Planck equation
\begin{equation}
    \boxed{\frac{dp}{dt} = \frac{p + vp'}{\tau} + \frac{A}{2}p''}
\end{equation}
\subsection{The velocity at time $t$}
the Langevin equation is just a inhomogenous ODE. The homogenouse solution for $\xi = 0$ is:
\begin{equation}
    v_h = e^{-t/\tau}
\end{equation}
We can find the private solution by variation of parameters. Setting:
\begin{equation}
    v_p = C(t)e^{-t/\tau}
\end{equation}
And plugging this into the equation, we get:
\begin{equation}
    \left(\dot{C} - \frac{1}{\tau}\right) e^{-t/\tau} = -\frac{C e^{-t/\tau}}{\tau} + \xi
\end{equation}
\begin{equation}
    \dot{C} = \xi e^{t/\tau}
\end{equation}
\begin{equation}
    c = \int_0^t \xi(t') e^{t'/\tau}dt'
\end{equation}
In total, we get:
\begin{equation}
    v(t) = v_0 e^{-t/\tau} + \int_0^t e^{-(t - t')/\tau} \xi(t')dt'
\end{equation}
Now, since $\xi$ is a gaussian, and the integral is a linear operator, the final $v(t)$ is also a gaussian. To find it, we need to only find the first and second moments. To do so, we will use the linearity of moments and the known moments of $\xi$:
\begin{equation}
    \mu_\xi = \langle \xi \rangle = 0 \qquad
    \sigma_\xi^2 = \langle \xi^2 \rangle - \langle \xi \rangle^2 = A
\end{equation}
\newline
The first moment:
\begin{equation}
\begin{gathered}
    \mu_v = \langle v_0 e^{-t/\tau} \rangle + \int_0^t e^{-(t - t')/\tau} \langle \xi(t') \rangle dt' = v_0 e^{-t/\tau} \\
    \sigma_v^2 = \langle v^2 \rangle - \langle v \rangle^2 = \\ =
    v_0^2 e^{-2t/\tau} + \int_0^t e^{-(t - t')/\tau} \langle v_0 e^{-t/\tau}\xi(t') \rangle dt' + \\ + \int_0^t e^{-2(t - t')/\tau} \langle \xi(t')\xi(t') \rangle dt' - v_0^2 e^{-2t/\tau} = \\ =
    v_0 e^{-t/\tau} \int_0^t e^{-(t - t')/\tau} \langle \xi(t') \rangle dt' + A \int_0^t e^{-2(t - t')/\tau} dt' =
    \frac{A \tau}{2} e^{2(t' - t)/\tau}\Big|_{t' = 0}^t  = \\ =
    \frac{A \tau}{2} \left(1 - e^{-2t/\tau}\right)
\end{gathered}
\end{equation}
So, this gives the velocity:
\begin{equation}
    p(v; t) = \frac{1}{\sqrt{2\pi}\frac{A \tau}{2} \left(1 - e^{-2t/\tau}\right)}\exp[-\frac{1}{2}\left(\frac{v - v_0 e^{-t/\tau}}{\frac{A \tau}{2} \left(1 - e^{-2t/\tau}\right)}\right)^2]
\end{equation}
\begin{equation}
    \boxed{p(v; t) = \frac{\sqrt{2}}{\sqrt{\pi}A \tau\left(1 - e^{-2t/\tau}\right)}\exp[-2\left(\frac{v - v_0 e^{-t/\tau}}{A \tau \left(1 - e^{-2t/\tau}\right)}\right)^2]}
\end{equation}
It's clear to see that at $t \rightarrow 0$ the distribution is a delta function around $v = v_0$, as expected.
\subsection{The $t \rightarrow \infty$ limit}
When taking $t \rightarrow \infty$, we get that all the exponents vanish:
\begin{equation}
    p(v; t \rightarrow \infty) = \frac{\sqrt{2}}{\sqrt{\pi}A \tau}\exp[-2\left(\frac{v}{A \tau}\right)^2]
\end{equation}
Which can be understood as the system forgetting the starting conditions and evolving to a steady state gaussian aroudnd $v = 0$.
\subsection{Thermodynamical interpertation}
Considering the steady state as the equillibrium of the system, we can expect the final distribution to be Boltzman distribution:
\begin{equation}
    p \propto e^{-H/kT}
\end{equation}
Comparing this to the final state we got and using the classical kinetic energy:
\begin{equation}
    2\left(\frac{v}{A \tau}\right)^2 = \frac{mv^2}{2kT}
\end{equation}
Which gives:
\begin{equation}
    \boxed{A = \sqrt{\frac{4kT}{m\tau^2} }}
\end{equation}
This means that the diffusion rate is related linearly to the temperature.
\end{document}